\apendice{Especificación de Requisitos}

\section{Introducción}




\section{Objetivos generales}

\section{Catálogo de requisitos}

\textbf{Requisitos funcionales}
\begin{itemize}
\tightlist
	\item RF-1 Inserción de documentos: el sistema debe permitir cargar documentos de licitaciones con el objetivo de construir y actualizar la base de datos vectorial. 
	\item RF.2 Procesado de los documentos: el sistema debe procesar de manera automática los documentos introducidos para prepararlos para que el RAG pueda usarlos.
	\begin{itemize}
	\tightlist
		\item RF-2.1 Extracción de información: el sistema debe extraer automáticamente el contenido textual de los documentos introducidos, ignorando páginas o documentos sin texto extraíble.
		\item RF-2.2 Conversión del contenido: el sistema debe transformar el texto extraído a un formato homogéneo, concretamente a \textit{markdown} para que sea más fácil su procesamiento.
		\begin{itemize}
		\tightlist
			\item RF-2.2.1 Conservar estructura básica: el sistema debe conservar la estructura básica del texto siempre que sea posible (secciones, títulos, saltos de línea, párrafos, etc).
		\end{itemize}
		\item RF-2.3 Limpieza y normalización del texto: el sistema debe limpiar y normalizar el texto que se extrae de los documentos.
		\begin{itemize}
		\tightlist
			\item RF-2.3.1 Eliminar caracteres: el sistema debe eliminar los caracteres irrelevantes o problemáticos.
			\item RF-2.3.2 Normalización: el sistema debe normalizar los espacios, saltos de línea y formatos inconsistentes.
		\end{itemize}
	\end{itemize}
	\item RF-3 Segmentación en chunks: el sistema debe dividir los documentos ya procesados en fragmentos de texto (\textit{chunks}) adecuados para el modelo de \textit{embeddings}.
	\begin{itemize}
	\tightlist
		\item RF-3.1 Tamaño: el sistema debe segmentar el texto respetando un máximo de tamaño por \textit{token} para evitar errores.
		\item RF-3.2 Solapamiento: el sistema debe aplicar solapamiento entre \textit{chunks} para preservar el contexto.
		\item RF-3.3 Contexto: cada \textit{chunk} debe tener una referencia al documento original y a su posición dentro de él.
	\end{itemize}
	\item RF-4 Generación de \textit{embeddings}: el sistema debe generar representaciones vectoriales (\textit{embeddings}) para cada \textit{chunk}.
	\item RF-5 Persistencia en base de datos vectorial: el sistema debe almacenar los \textit{chunks} y sus \textit{embeddings} en una base de datos vectorial.
	\begin{itemize}
	\tightlist
		\item RF-5.1 Almacenar chunks: el sistema debe almacenar el contenido textual de cada \textit{chunk}.
		\item RF-5.2 Almacenar embeddings: el sistema debe almacenar el \textit{embedding} asociado a cada \textit{chunk}.
		\item RF-5.3 Almacenar metadatos: el sistema debe almacenar metadatos relevantes como el nombre del archivo, el titulo y el segmento incluido en el chunk.
		\item RF-5.4 Actualización: el sistema debe permitir la actualización o recreación de la colección vectorial en el caso de incluir documentos nuevos.
	\end{itemize}
	\item RF-6 Recuperación de información por similitud: el sistema debe recuperar información relevante a partir de las consultas de los usuarios.
	\begin{itemize}
	\tightlist
		\item RF-6.1 Generación de \textit{embedding}: el sistema debe generar el \textit{embedding} de la consulta del usuario.
		\item RF-6.2 Búsqueda por similitud: el sistema debe realizar una búsqueda por similitud en la base de datos vectorial.
	\end{itemize}
	\item RF-7 Construcción del \textit{prompt} con contexto: el sistema debe construir un \textit{prompt}, para alimentar al LLM, que incluya los \textit{chunks} recuperados, la pregunta del usuario y reglas de uso del contexto.
	\item RF-8 Generación de respuestas mediante un LLM: el sistema debe enviar el \textit{prompt} generado a un modelo LLM y mostrar la respuesta de dicho modelo al usuario.
	\begin{itemize}
	\tightlist
		\item RF-8.1 Evitar alucinaciones: El sistema no debe responder cuando no haya contexto relevante suficiente.
	\end{itemize}
	\item RF-9 Trazabilidad de respuesta: el sistema debe devolver junto con la respuesta los metadatos del documento origen y los \textit{chunks} empleados para la generación de la respuesta. 
	\item RF-10 Gestión de usuarios: el sistema debe permitir la gestión de usuarios a través de la aplicación web.
	\begin{itemize}
	\tightlist
		\item RF-10.1 Registros y autentificación: el sistema debe permitir el registro y autentificación de los usuarios.
		\item RF-10.2 Sesiones: el sistema debe gestionar las sesiones de los usuarios de forma segura.
		\item RF-10.3 Roles: el sistema debe permitir asignar a los usuarios el rol de administradores o usuarios normales.
		\item RF.10.4 Administradores: los administradores deben poder añadir usuarios, eliminarlos y cambiar su rol.
	\end{itemize}
	\item RF-11 Interfaz de consultas: los usuarios autentificados podrán realizar consultas y ver las respuestas del sistema.
	\begin{itemize}
	\tightlist
		\item RF-11.1 Historial de consultas: los usuarios podrán ver su historial de consultas.
	\end{itemize}
\end{itemize}

\textbf{Requisitos no funcionales}
\begin{itemize}
\tightlist
	\item RNF-1 Reproducibilidad: el proyecto debe ejecutarse en un entorno reproducible.
	\item RNF-2 Rendimiento: el sistema debe ser capaz de procesar los documentos sin intervención manual. Además, el sistema tiene que responder en un tiempo razonable.
	\item RNF-3 Robustez ante errores: el sistema debe continuar ejecutándose ante fallos esperables en datos, red o servicios externos.
	\begin{itemize}
	\tightlist
		\item RNF-3.1 Documentos corruptos: el sistema debe saltarse los documentos corruptos o sin texto extraible y continuar con el resto de documentos.
		\item RNF-3.2 \textit{Timeouts}: el sistema debe manejar los \textit{timeouts} del LLM devolviendo mensajes controlados al usuario.
	\end{itemize}
	\item RNF-4 Trazabilidad técnica: el sistema debe realizar \textit{logs} de carga de modelo, lectura de PDFs, \textit{chunking}, \textit{embeddings}, guardado y consulta.
	\item RNF-5 Seguridad: el sistema debe gestionar de forma segura la información sensible de los usuarios y el acceso a la aplicación.
	\begin{itemize}
	\tightlist
		\item RNF-5.1 Contraseñas: las contraseñas deben almacenarse \textit{hasheadas}.
		\item RNF-5.2 Control de acceso por roles: las funciones administrativas deben estar protegidas por rol para que solo las pueda realizar el administrador.
		\item RNF-5.3 Protección de secretos: las claves sensibles no deben estar públicas en el repositorio.
	\end{itemize}
	\item RNF-6 Usabilidad: la interfaz debe permitir a los usuarios realizar las consultas y ver el historial de manera sencilla. Además, los textos deben ser legibles.
	\item RNF-7 Calidad del \textit{software}: debe existir un conjunto de pruebas que asegure el correcto funcionamiento del programa. 
	\item RNF-8 Disponibilidad: la web debe garantizar un nivel de disponibilidad adecuado que permita a los usuarios acceder a la aplicación y realizar consultas.
	\item RNF-9 Recuperación y tolerancia a fallos: el sistema debe disponer de mecanismos que permitan la recuperación de información y la restauración del servicio tras fallos, garantizando la integridad de los datos.
\end{itemize}

\section{Especificación de requisitos}

\subsection{Actores}
Los actores del sistema serán:
\begin{itemize}
\tightlist
	\item \textbf{Usuario anónimo}: usuario que no ha iniciado sesión (solo puede registrarse y autentificarse)
	\item \textbf{Usuario autentificado}: usuario que puede acceder a la aplicación sus acciones en al aplicación cambiarán dependiendo de su rol. 
	\begin{itemize}
	\tightlist
	 	\item \textbf{Usuario normal}: usuario autentificado que puede consultar el RAG y ver su historial
	 	\item \textbf{Usuario administrador}: usuario autentificado con permisos de gestión de usuarios y de carga y actualización documental
	 \end{itemize}
	 \item \textbf{Sistema}: es el encargado del procesamiento, \textit{chunking}, \textit{embeddings}, guardado y consulta. 
\end{itemize}

\subsection{Casos de uso}
Los casos de uso definidos para el proyecto son:
\begin{itemize}
\tightlist
	\item CU-1 Registro de usuarios
	\item CU-2 Iniciar sesión
	\item CU-3 Cerrar sesión
	\item CU-4 Realizar consulta al RAG
	\item CU-5 Consultar historial de consultas
	\item CU-6 Cargar documentos de licitaciones
	\item CU-7 Procesar documentos cargados
	\item CU-8 Indexar documentos
	\item CU-9 Actualizar base de datos vectorial
	\item CU-10 Administrar usuarios
\end{itemize}

\begin{table}[p]
	\centering
	\begin{tabularx}{\linewidth}{ p{0.21\columnwidth} p{0.71\columnwidth} }
		\toprule
		\textbf{CU-1}    & \textbf{Registro de usuarios}\\
		\toprule
		\textbf{Versión}              & 1.0    \\
		\textbf{Autor}                & Lydia Blanco Ruiz \\
		\textbf{Requisitos asociados} & RF-10.1 \\
		\textbf{Descripción}          & Permite crear una cuenta en la aplicación web. \\
		\textbf{Precondición}         & Que el usuario no este autentificado. \\
		\textbf{Acciones}             &
		\begin{enumerate}
			\def\labelenumi{\arabic{enumi}.}
			\tightlist
			\item El usuario accede a ``Registro''.
			\item El usuario introduce el nombre, email y contraseña.
			\item El sistema valida el formato y que el email no este duplicado.
			\item El sistema almacena el usuario con la contraseña \textit{hasheada}.
			\item El sistema confirma el registro.
			\item El sistema crea la sesión y redirige a la página de consultas.
		\end{enumerate}\\
		\textbf{Postcondición}        & Cuenta creada y usuario autentificado \\
		\textbf{Excepciones}          & 
		\begin{enumerate}
			\def\labelenumi{\arabic{enumi}.}
			\tightlist
			\item Tiene algún campo vacío. Muestra un error con el mensaje ``Completa este campo''.
			\item El email ya existe. Muestra un error con el mensaje ``Ya existe un usuario con ese email''.
			\item El email no cumple la validación. Muestra un error con el mensaje ``Invalid email address''.
			\item La contraseña no cumple la política. Muestra un error con el mensaje ``Aumenta la longitud del texto a 6 caracteres como mínimo''.
			\item Las contraseñas no coinciden. Muestra un error con el mensaje ``Las contraseñas no coinciden''.
			\item El nombre tiene menos de 2 caracteres. Muestra un error con el mensaje ``Aumenta la longitud del texto a 2 caracteres o más''.
		\end{enumerate}\\
		\textbf{Importancia}          & Alta  \\
		\bottomrule
	\end{tabularx}
	\caption{CU-1 Registro de usuarios.}
\end{table}

\begin{table}[p]
	\centering
	\begin{tabularx}{\linewidth}{ p{0.21\columnwidth} p{0.71\columnwidth} }
		\toprule
		\textbf{CU-2}    & \textbf{Iniciar sesión}\\
		\toprule
		\textbf{Versión}              & 1.0    \\
		\textbf{Autor}                & Lydia Blanco Ruiz \\
		\textbf{Requisitos asociados} & RF-10.1, RF-10.2 \\
		\textbf{Descripción}          & Autentifica al usuario y abre una sesión segura. \\
		\textbf{Precondición}         & Que el usuario este registrado y no este autentificado. \\
		\textbf{Acciones}             &
		\begin{enumerate}
			\def\labelenumi{\arabic{enumi}.}
			\tightlist
			\item El usuario accede a ``Iniciar sesión''.
			\item El usuario introduce el email y la contraseña.
			\item El sistema verifica los credenciales.
			\item El sistema crea la sesión y redirige a la página de consultas.
		\end{enumerate}\\
		\textbf{Postcondición}        & Sesión activa para el usuario. \\
		\textbf{Excepciones}          & 
		\begin{enumerate}
			\def\labelenumi{\arabic{enumi}.}
			\tightlist
			\item Tiene algún campo vacío. Muestra un error con el mensaje ``Completa este campo''.
			\item El email o la contraseña no coinciden con ningún usuario registrado. Muestra un error con el mensaje ``Email o contraseña incorrectos.''.
			\item El email no cumple la validación. Muestra un error con el mensaje ``Invalid email address''.
			\item La contraseña no cumple la política. Muestra un error con el mensaje ``Aumenta la longitud del texto a 6 caracteres como mínimo''.		
		\end{enumerate}\\
		\textbf{Importancia}          & Alta  \\
		\bottomrule
	\end{tabularx}
	\caption{CU-2 Iniciar sesión.}
\end{table}

\begin{table}[p]
	\centering
	\begin{tabularx}{\linewidth}{ p{0.21\columnwidth} p{0.71\columnwidth} }
		\toprule
		\textbf{CU-3}    & \textbf{Cerrar sesión}\\
		\toprule
		\textbf{Versión}              & 1.0    \\
		\textbf{Autor}                & Lydia Blanco Ruiz \\
		\textbf{Requisitos asociados} & RF-10.2 \\
		\textbf{Descripción}          & Finaliza la sesión del usuario. \\
		\textbf{Precondición}         & Que el usuario este autentificado. \\
		\textbf{Acciones}             &
		\begin{enumerate}
			\def\labelenumi{\arabic{enumi}.}
			\tightlist
			\item El usuario pulsa el botón ``Cerrar sesión''.
			\item El sistema invalida la sesión.
			\item El sistema redirige a la primera de inicio.
		\end{enumerate}\\
		\textbf{Postcondición}        & Sesión finalizada. \\
		\textbf{Excepciones}          & 
		\begin{enumerate}
			\def\labelenumi{\arabic{enumi}.}
			\tightlist
			\item Error al invalidar sesión. Muestra un mensaje de error controlado. 	
		\end{enumerate}\\
		\textbf{Importancia}          & Media  \\
		\bottomrule
	\end{tabularx}
	\caption{CU-3 Cerrar sesión.}
\end{table}

\begin{table}[p]
	\centering
	\begin{tabularx}{\linewidth}{ p{0.21\columnwidth} p{0.71\columnwidth} }
		\toprule
		\textbf{CU-4}    & \textbf{Realizar consulta al RAG}\\
		\toprule
		\textbf{Versión}              & 1.0    \\
		\textbf{Autor}                & Lydia Blanco Ruiz \\
		\textbf{Requisitos asociados} & RF-6, RF-7, RF-8, RF-8.1, RF-9, RF-11 \\
		\textbf{Descripción}          & El usuario hace una pregunta, el sistema recupera contexto por similitud, construye el prompt y genera respuesta que incluya las fuentes. \\
		\textbf{Precondición}         & Que el usuario este autentificado y la base de datos vectorial este disponible. \\
		\textbf{Acciones}             &
		\begin{enumerate}
			\def\labelenumi{\arabic{enumi}.}
			\tightlist
			\item El usuario escribe y envía una pregunta.
			\item El sistema genera \textit{embeddings} de la consulta.
			\item El sistema busca por similitud en la base de datos vectorial.
			\item El sistema construye un \textit{prompt} con los \textit{chunks} y las reglas de uso del contexto.
			\item El sistema llama al LLM y genera respuestas.
			\item El sistema llama al LLM y genera una respuesta.
			\item El sistema adjunta la trazabilidad de la respuesta, la cual incluye los documentos, metadatos y \textit{chunks} utilizados.
			\item El sistema guarda la consulta en el historial del consultas del usuario.
		\end{enumerate}\\
		\textbf{Postcondición}        & Respuesta mostrada y consulta registrada. \\
		\textbf{Excepciones}          & 
		\begin{enumerate}
			\def\labelenumi{\arabic{enumi}.}
			\tightlist
			\item Cuando no hay contexto suficiente. El sistema no responde a la pregunta (para evitar alucinaciones) y muestra un aviso de que falta contexto.
			\item Timeout del LLM. El sistema muestra un mensaje controlado. 	
		\end{enumerate}\\
		\textbf{Importancia}          & Alta  \\
		\bottomrule
	\end{tabularx}
	\caption{CU-4 Realizar consulta al RAG.}
\end{table}